% !TEX root = main.tex
\chapter*{Abstract\label{Abstract}}
\addcontentsline{toc}{chapter}{Abstract}
\fancyhead[RE]{\slshape \textsf{Abstract}} %
\fancyhead[LO]{\slshape \textsf{Abstract}} %

\begin{otherlanguage*}{english}

Earth-observation satellites acquire increasing volumes of multispectral data
under constraints on downlink capacity, memory, and on-board processing. This
thesis implements the compression and decompression pipeline of SORTENY, a
neural multispectral image codec, in C and extends it with spatial quality
control that does not require manual selection of regions of interest.

Quality is assigned block-wise through maps generated automatically from the
spectral and spatial content of the image. Detectors identify priority blocks,
and quality policies determine how fidelity is distributed between those
blocks and the background. Two complementary strategies were studied: one
prioritizes bit savings by degrading non-priority regions, whereas the other
assigns a fixed high quality to the detected region to favor its preservation.

The evaluation included 120 eight-band Sentinel-2 crops. The coded bitrate was
subsequently measured with the reference pipeline on 20 selected crops. The
bit-saving strategy reduced the mean bitrate by 8.84\% relative to applying
uniform quality to the whole image, while maintaining or slightly improving
fidelity in the prioritized regions and concentrating degradation in the
background. The preservation strategy protected those regions further at the
cost of a smaller bitrate reduction.

The C implementation was also evaluated on a Raspberry Pi 3B+ as a
representative resource-constrained platform. Relative to the initial C version
running with the same four threads, the introduced optimizations reduced mean
total runtime by 17.0\%, equivalent to a $1.205\times$ speed-up, with no
detected functional changes or relevant increase in memory use. Automatic
quality-map generation accounted for less than 0.05\% of the codec runtime.
Therefore, spatial decision has negligible cost, and neural codec execution
remains the main bottleneck. The current C bitstream still lacks the final
entropy or arithmetic coder. The Raspberry is used as a computational proxy
for a resource-constrained on-board platform; validation on the final hardware
of a specific mission remains outside the scope of this work.

\textbf{Keywords:} multispectral image compression, learned compression,
fixed-quality compression, region of interest, on-board processing, embedded
systems, Raspberry Pi, quality map.

\end{otherlanguage*}
